%% =====================================================================
%%  T-10-03  Limited Numbers and Deadlines
%%  -------------------------------------------------------------------
%%  One idea per file. Core is mandatory. Slots in canonical order.
%%  Max 700 words. No layout commands. No filler. New source material
%%  for this entry goes in sources/<BOOK>/T-10-03.tex -- never by
%%  rewriting this file (docs/RULES.md R0.1).
%% =====================================================================
%% @kind: topic
%% @id: T-10-03
%% @title: Limited Numbers and Deadlines
%% @subtitle: The two forms scarcity is delivered in
%% @chapter: c10
%% @order: 30
%% @status: stable
%% @sources: B4
%% @updated: 2026-09-19

\topic{T-10-03}{Limited Numbers and Deadlines}{The two forms scarcity is delivered in}

\begin{core}
Scarcity reaches a person in two standard forms: a limited quantity and a closing date. Both are announcements about the world that can be made without changing anything in it, and both work hardest when they arrive together and suddenly.
\end{core}
\begin{mechanism}
A limited quantity works through competition. The number is not only a constraint but a signal that others want the thing, which imports crowd evidence into the scarcity and doubles its force --- the object is rarer and independently confirmed as desirable.

A deadline works through the loss of the option to decide later. Postponement is the ordinary answer to a difficult judgement, because more information usually arrives; a closing date removes that answer and forces a decision under exactly the conditions of ambiguity in which shortcuts are followed most readily. The deadline does not make the thing more valuable. It makes the decision worse, and a worse decision is an easier one to steer.

Novelty strengthens both. A scarcity that has just appeared is more alarming than one that has always been there, because the change implies information you do not have. The announcement therefore matters more than the condition.
\end{mechanism}
\begin{tells}
  \item A quantity or a date is stated with precision and no explanation.
  \item Both arrive together.
  \item The scarcity is new, and its newness is emphasised.
  \item You are deciding under time pressure about something you would normally take a day over.
  \item Others are visibly competing for the same limited thing.
  \item The announcement serves the seller more than it informs you.
\end{tells}
\begin{moves}
  \item Announce the limit before the request, with a specific number or date rather than a general warning.
  \item Present the scarcity as recent. A change implies information and is more alarming than a standing condition.
  \item Combine quantity and deadline, and let each confirm the other.
  \item Make the competition visible; the number works harder when other claimants can be seen.
\end{moves}
\begin{counters}
  \item Ask when the limit was set and by whom. A limit the seller chose is an instrument, not a fact.
  \item Check the history. If the same deadline has passed before without consequence, it is not a deadline.
  \item Separate the two claims: is the thing scarce, and is it desirable? Answer each without reference to the other.
  \item Remove the time pressure and see whether the judgement survives. A decision that only works inside the window was never about the object.
  \item Ask for the terms in writing, dated, with the limit stated. Manufactured limits are usually unwilling to be written down.
\end{counters}
\begin{cost}
Both forms are cheap to falsify and therefore cheap to detect. A deadline that passes and returns, or a final unit that is followed by another, tells the customer everything about the seller at once --- and the discovery is more damaging than the extra sale was worth. Used honestly, both are simply information about real constraints, and the honest version has most of the effect at none of the cost, which is the strongest argument for it.
\end{cost}
\begin{field}
One room left, nine people viewing it, a rate expiring at midnight: three instruments, not three facts.
\end{field}

\sources{B4}
\seesources
\seealso{T-10-01, T-10-02, T-24-03, T-15-01}
